\documentclass[12 pt, letterpaper]{hmcpset}
\usepackage{hw-preamble}

\name{Jayden Thadani}
\class{MATH 145 --- Topics in Geometry and Topology}
\assignment{Homework 10}
\duedate{16th April 2025}

\begin{document}

Recall that $E = \{ \O \}$ denotes the empty simplicial complex. Recall also that $C$ and $S$ denote the cone and suspension operations.

\begin{problem}[52.]
	\begin{enumerate}
		\item Sketch the simplicial complexes $CE$, $CCE$,  $C C C E$, $C C C C E$.

		\item Determine the dimension of $C C C C C E$ and count the top-dimensional simplices.
	\end{enumerate}
\end{problem}

\pagebreak

\begin{problem}[53.]
	\begin{enumerate}
		\item Sketch the simplicial complexes $SE$, $SSE$, $SSSE$.

		\item Determine the dimension of $SSSSE$ and count the top-dimensional simplices.
	\end{enumerate}
\end{problem}

\pagebreak


\begin{problem}[54.]
	Given a poset $P$, we define two related simplicial complexes as follows. In both cases, $P$ is the vertex set. The \emph{simplicial complex of chains} $\operatorname{Ch} P$ is defined
	\[
		\sigma = \{ x_{0}, x_{1}, \dots, x_{n} \} \in \operatorname{Ch} P
	\]
	if and only if every pair of elements $x_{i}, x_{j}$ is comparable.

	The \emph{simplicial compelx of antichains} $\operatorname{Anti} P$ is defined
	\[
		\sigma = \{ x_{0}, x_{1}, \dots, x_{n} \} \in \operatorname{Anti} P
	\]
	if and only if no pair of distinct elements $x_{i}, x_{j}$ is comparable.

	Let $P$ be the poset defined by the following \emph{Hasse diagram}:

	In other words, $a$ is greater than everything else, $e$ is less than everything else, $b, c, d$ are pairwise incomparable with each other.

	\begin{enumerate}
		\item Describe and sketch $\operatorname{Ch} P$.

		\item Describe and sketch $\operatorname{Anti} P$.
	\end{enumerate}
\end{problem}

\pagebreak

\begin{problem}[55.]
	\begin{enumerate}
		\item Draw a regular hexagon whose vertices lie on the unit circle in the plane, and calculate the distances which occur between pairs of these vertices.

		\item For all values of $R \ge 0$, identify the Vietoris-Rips complex $\operatorname{VR}(X, R)$ where $X$ is the set of vertices of a regular hexagon.
	\end{enumerate}
\end{problem}

\pagebreak

\begin{problem}[56.]
	\begin{enumerate}
		\item Exhibit a simplicial complex $P$ with $\mathrm{b}_{0}(P) = 4$ and $\chi(P) = 7$.

		\item Exhibit a simplicial complex $P$ with $\mathrm{b}_{0}(P) = 7$ and $\chi(P) = 4$. 
	\end{enumerate}
\end{problem}

\pagebreak

\begin{problem}[57.]
	Let $X$ be a simplicial complex $X$. Removing the empty simplex, we consider $X \setminus \{ \O \}$ as a poset with $\sigma \le \tau$ if $\sigma$ is a face of $\tau$. We define the \emph{derived complex of $X$}
	\[
		 X^{(1)} = \operatorname{Ch} (X \setminus \{ \O \})
	\]
	to be the simplicial complex of chains of $X \setminus \{ \O \}$.

	In each part, sketch the derived complex of the given complex. Try to find a good way to draw it that reveals what is actually going on with this construction.

	\begin{enumerate}
		\item $X = B_{0}$.

		\item $X = B_{1}$.

		\item $X = \partial B_{2}$.

		\item $X = B_{2}$.

		\item $X = $
	\end{enumerate}
\end{problem}

\pagebreak

\begin{problem}[58.]
	Consider the Whitehead equivalence relation, generated by elementary collapses. Two of the four complexes drawn below are Whitehead equivalent to each other, as are the other two.

	For each equivalent pair, demonstrate the Whitehead equivalence by exhibiting a sequence of elementary collapses and expansions that lead from one to the other.
\end{problem}

\end{document}
